\documentclass[UTF8,zihao=-4]{ctexart}
\usepackage{geometry}
\geometry{a4paper,margin=2.05cm,headheight=15pt}
\usepackage{amsmath,amssymb,amsthm,mathtools}
\usepackage{array,booktabs,longtable,tabularx,multirow}
\usepackage{enumitem}
\usepackage{multicol}
\usepackage{xcolor}
\usepackage[most]{tcolorbox}
\usepackage{fancyhdr}
\usepackage{hyperref}
\usepackage{tikz}
\usepackage{eso-pic}
\usetikzlibrary{arrows.meta,positioning,calc}
\usepackage{lastpage}
\usepackage{fontspec}

\setmainfont{DejaVu Serif}
\setsansfont{DejaVu Sans}
\setmonofont{DejaVu Sans Mono}
\setCJKmainfont{Noto Serif CJK SC}
\setCJKsansfont{Noto Sans CJK SC}
\setCJKmonofont{Noto Sans Mono CJK SC}

\definecolor{main}{HTML}{1F4E79}
\definecolor{accent}{HTML}{B45F06}
\definecolor{softblue}{HTML}{EAF3F8}
\definecolor{softorange}{HTML}{FFF3E8}
\definecolor{softgreen}{HTML}{EAF7EA}
\definecolor{softred}{HTML}{FBEAEA}
\definecolor{softgray}{HTML}{F5F5F5}

\hypersetup{colorlinks=true,linkcolor=main,urlcolor=main,citecolor=main,pdftitle={离散期末复习},pdfauthor={ChatGPT}}

\pagestyle{fancy}
\fancyhf{}
\lhead{离散期末复习}
\rhead{数理逻辑 · 集合论 · 图论}
\cfoot{\thepage/\pageref{LastPage}}

\setlist[itemize]{leftmargin=2em,itemsep=0.18em,topsep=0.3em}
\setlist[enumerate]{leftmargin=2.2em,itemsep=0.22em,topsep=0.3em}
\renewcommand{\arraystretch}{1.22}

\newcommand{\kw}[1]{\textbf{\color{main}#1}}
\newcommand{\warn}[1]{\textbf{\color{accent}#1}}
\newcommand{\set}[1]{\left\{#1\right\}}
\newcommand{\pow}{\mathcal P}
\newcommand{\dom}{\operatorname{dom}}
\newcommand{\ran}{\operatorname{ran}}
\newcommand{\id}{\operatorname{id}}
\newcommand{\True}{\mathrm T}
\newcommand{\False}{\mathrm F}
\newcommand{\N}{\mathbb N}
\newcommand{\Z}{\mathbb Z}
\newcommand{\I}{I_A}

\newtcolorbox{corebox}[1]{enhanced,breakable,colback=softblue,colframe=main!75!black,title={#1},fonttitle=\bfseries,arc=1.2mm,boxrule=0.8pt,left=1.3mm,right=1.3mm,top=1mm,bottom=1mm}
\newtcolorbox{methodbox}[1]{enhanced,breakable,colback=softgreen,colframe=green!45!black,title={#1},fonttitle=\bfseries,arc=1.2mm,boxrule=0.8pt,left=1.3mm,right=1.3mm,top=1mm,bottom=1mm}
\newtcolorbox{pitfallbox}[1]{enhanced,breakable,colback=softred,colframe=red!55!black,title={#1},fonttitle=\bfseries,arc=1.2mm,boxrule=0.8pt,left=1.3mm,right=1.3mm,top=1mm,bottom=1mm}
\newtcolorbox{exambox}[1]{enhanced,breakable,colback=softorange,colframe=accent!85!black,title={#1},fonttitle=\bfseries,arc=1.2mm,boxrule=0.8pt,left=1.3mm,right=1.3mm,top=1mm,bottom=1mm}


\definecolor{watermarkgray}{HTML}{777777}
\newcommand{\WatermarkText}{Jerry 制作｜仅供个人复习使用｜禁止盗印与商用传播}
\AddToShipoutPictureBG{%
  \begin{tikzpicture}[remember picture,overlay]
    \node[rotate=35,scale=1.35,text=watermarkgray,opacity=0.14] at (current page.center) {Jerry 制作｜仅供个人复习使用｜禁止盗印与商用传播};
  \end{tikzpicture}%
}


\begin{document}

\begin{titlepage}
\centering
\vspace*{1.1cm}
{\zihao{1}\bfseries 离散期末复习\par}
\vspace{0.35cm}
{\zihao{3}\color{main} 根据老师考点 + 教材标黄内容修订\par}
\vspace{0.9cm}
\begin{tcolorbox}[enhanced,colback=softblue,colframe=main,width=0.86\textwidth,arc=2mm,boxrule=0.9pt]
\centering
\large
期末分值：\kw{数理逻辑 30 分}，\kw{集合论 34 分}，\kw{图论 36 分}。\\[0.35em]
本版已按你的补充调整：\kw{命题推理 2.8 放在一阶逻辑前}；图论拆为\kw{图的基本概念与矩阵表示}、\kw{特殊的图}、\kw{树}三大部分。
\end{tcolorbox}
\vspace{0.8cm}
\begin{tabular}{>{\centering\arraybackslash}p{3.3cm}>{\centering\arraybackslash}p{2.4cm}p{8.2cm}}
\toprule
模块 & 分值 & 复习重点 \\
\midrule
数理逻辑 & 30 & 命题符号化、公式真值、公式类型、等值式、主范式、构造证明、一阶逻辑符号化、前束范式 \\
集合论 & 34 & 集合运算与证明、关系运算、关系性质与闭包、等价关系、偏序关系、函数、特征函数、自然映射、复合与反函数 \\
图论 & 36 & 握手定理、连通性、割集、矩阵表示、可达矩阵、最短路、图着色、二部图、欧拉图、哈密顿图、树、Kruskal、Huffman、二叉树遍历 \\
\bottomrule
\end{tabular}
\vfill
{\large 建议：公式先背，判定模板再练，计算题最后做速度训练。}
\end{titlepage}

\tableofcontents
\newpage

\section{考试结构与总复习路线}

\begin{corebox}{老师给出的总考点}
\begin{enumerate}
  \item \kw{数理逻辑部分}：命题符号化（包括一阶逻辑）、求公式的值（包括一阶逻辑）、判断公式的类型（包括一阶逻辑）、判断等值式（包括一阶逻辑）、求主范式、求前束范式、推理（构造证明法）。
  \item \kw{集合论部分}：集合运算、证明集合的包含或相等、关系的运算、关系的性质与闭包、偏序关系、等价关系、函数的定义与性质、特殊函数（特征函数、自然映射）、复合函数与反函数。
  \item \kw{图论部分}：握手定理的应用、判断有向图的连通性、求点割集和边割集、图的矩阵表示与可达矩阵、最短路径算法、图着色、二部图、欧拉图、哈密顿图、最小生成树算法 Kruskal、哈夫曼算法（包括编码）、二叉树遍历（包括波兰符号和逆波兰符号表达式）。
\end{enumerate}
\end{corebox}

\begin{methodbox}{复习顺序}
\begin{enumerate}
  \item 先背定义和公式：逻辑等值式、集合运算律、关系性质、图论判定定理。
  \item 再练计算：主范式、前束范式、关系闭包、矩阵幂、可达矩阵、最短路、Kruskal、Huffman。
  \item 最后练证明：集合相等、推理证明、等价关系、偏序关系、二部图/欧拉图/哈密顿图判定。
\end{enumerate}
\end{methodbox}

\section{数理逻辑部分（30 分）}

\subsection{命题逻辑基础与符号化}

\begin{corebox}{命题与联结词}
命题是能判断真假的陈述句。常用命题变元为 $p,q,r,\ldots$。
\[
\begin{array}{c|c|c}
\text{符号} & \text{读法} & \text{真值条件}\\\hline
\neg p & \text{非 }p & p\text{ 假时真}\\
p\land q & p\text{ 且 }q & p,q\text{ 都真时真}\\
p\lor q & p\text{ 或 }q & p,q\text{ 至少一个真时真}\\
p\to q & \text{若 }p\text{ 则 }q & \text{仅 }p=\True,q=\False\text{ 时假}\\
p\leftrightarrow q & p\text{ 当且仅当 }q & p,q\text{ 同真同假时真}
\end{array}
\]
\end{corebox}

\begin{methodbox}{自然语言符号化速查}
\begin{center}
\begin{tabularx}{\textwidth}{>{\bfseries}p{3.4cm}X}
\toprule
自然语言 & 公式 \\
\midrule
若 $p$，则 $q$ & $p\to q$ \\
只有 $p$，才 $q$ & $q\to p$ \\
$p$ 仅当 $q$ & $p\to q$ \\
$p$ 是 $q$ 的充分条件 & $p\to q$ \\
$p$ 是 $q$ 的必要条件 & $q\to p$ \\
$p$ 当且仅当 $q$ & $p\leftrightarrow q$ \\
$p$ 除非 $q$ & $\neg q\to p$，等价于 $p\lor q$ \\
$p,q$ 至少一个成立 & $p\lor q$ \\
$p,q$ 恰有一个成立 & $(p\lor q)\land\neg(p\land q)$ \\
并非 $p$ 且 $q$ & $\neg(p\land q)$，不是 $\neg p\land q$ \\
\bottomrule
\end{tabularx}
\end{center}
\end{methodbox}

\subsection{求公式的值与判断公式类型}

\begin{corebox}{四类公式}
设 $A$ 为命题公式。
\begin{itemize}
  \item \kw{重言式/永真式}：任意赋值下 $A$ 都为真。
  \item \kw{矛盾式/永假式}：任意赋值下 $A$ 都为假。
  \item \kw{可满足式}：至少存在一个赋值使 $A$ 为真。
  \item \kw{偶然式}：有时真，有时假。
\end{itemize}
\end{corebox}

\begin{methodbox}{判断公式类型}
\begin{enumerate}
  \item 变元少：列真值表。
  \item 结构明显：用等值式化简。
  \item 要证 $A$ 是重言式：化到 $1$；要证 $A$ 是矛盾式：化到 $0$。
\end{enumerate}
\end{methodbox}

\subsection{常用等值式与主范式（对应教材 2.5）}

\begin{corebox}{命题逻辑常用等值式}
\begin{multicols}{2}
\begin{align*}
&\neg\neg p\equiv p,\\
&p\land q\equiv q\land p,\\
&p\lor q\equiv q\lor p,\\
&(p\land q)\land r\equiv p\land(q\land r),\\
&(p\lor q)\lor r\equiv p\lor(q\lor r),\\
&p\land(q\lor r)\equiv(p\land q)\lor(p\land r),\\
&p\lor(q\land r)\equiv(p\lor q)\land(p\lor r),\\
&p\land p\equiv p,\quad p\lor p\equiv p,\\
&p\land 1\equiv p,\quad p\lor 0\equiv p,\\
&p\land 0\equiv 0,\quad p\lor 1\equiv 1,\\
&p\land\neg p\equiv 0,\quad p\lor\neg p\equiv 1,\\
&\neg(p\land q)\equiv\neg p\lor\neg q,\\
&\neg(p\lor q)\equiv\neg p\land\neg q,\\
&p\to q\equiv\neg p\lor q,\\
&p\leftrightarrow q\equiv(p\to q)\land(q\to p).
\end{align*}
\end{multicols}
\end{corebox}

\begin{corebox}{主析取范式与主合取范式}
给定变元顺序 $p_1,p_2,\ldots,p_n$：
\begin{itemize}
  \item \kw{小项} $m_i$：每个变元出现一次，用合取连接。小项只在一个赋值下为真。
  \item \kw{大项} $M_i$：每个变元出现一次，用析取连接。大项只在一个赋值下为假。
  \item \kw{主析取范式}：把公式为真的行对应的小项析取起来：$\bigvee m_i$。
  \item \kw{主合取范式}：把公式为假的行对应的大项合取起来：$\bigwedge M_i$。
\end{itemize}
\end{corebox}

\begin{methodbox}{主范式计算模板}
\begin{enumerate}
  \item 确定变元顺序，例如 $p,q,r$。
  \item 列真值表，求原公式每一行的真值。
  \item 真值为 $1$ 的行写成小项，得到主析取范式。
  \item 真值为 $0$ 的行写成大项，得到主合取范式。
\end{enumerate}
编号规则：通常把 $(0,0,0),(0,0,1),\ldots,(1,1,1)$ 看作二进制编号。小项中 $1$ 对应原变元，$0$ 对应否定；大项相反，$0$ 对应原变元，$1$ 对应否定。
\end{methodbox}

\subsection{推理与构造证明法（对应教材 2.8，放在一阶逻辑前）}

\begin{corebox}{命题逻辑常用推理规则}
\begin{center}
\begin{tabularx}{\textwidth}{p{3.1cm}p{5.2cm}X}
\toprule
规则 & 形式 & 说明 \\
\midrule
前提引入 & 已知前提可直接写入证明 & 证明序列第一类来源 \\
结论引入 & 已证结论可作为后续证明前提 & 用已得结论继续推 \\
置换规则 & $A\equiv B$ 时可互换 & 等值式替换 \\
假言推理 MP & $A\to B,\;A\vdash B$ & 最常用 \\
拒取式 MT & $A\to B,\;\neg B\vdash \neg A$ & 逆否推理 \\
假言三段论 & $A\to B,\;B\to C\vdash A\to C$ & 蕴含链 \\
析取三段论 & $A\lor B,\;\neg A\vdash B$ & 排除一支 \\
附加规则 & $A\vdash A\lor B$ & 加一个析取支 \\
化简规则 & $A\land B\vdash A$；$A\land B\vdash B$ & 从合取式取出一项 \\
合取引入 & $A,B\vdash A\land B$ & 两式合并 \\
构造性二难 & $(A\to B)\land(C\to D),\;A\lor C\vdash B\lor D$ & 分情况推出析取结论 \\
\bottomrule
\end{tabularx}
\end{center}
\end{corebox}

\begin{methodbox}{构造证明法格式}
证明 $P_1,P_2,\ldots,P_k\vdash Q$：
\begin{enumerate}
  \item 每一行写一个公式。
  \item 每一行右侧标明依据：前提引入、等值置换、假言推理、化简、合取引入等。
  \item 最后一行必须得到 $Q$。
\end{enumerate}
\end{methodbox}

\begin{exambox}{证明题思路}
若题目给出多个蕴含式，例如 $p\to q,q\to r$，优先尝试假言三段论得到 $p\to r$。若给出 $A\land B$，优先化简；若目标是合取式，先分别证明两边，再用合取引入。
\end{exambox}

\subsection{一阶逻辑：符号化与公式求值（对应教材 2.6）}

\begin{corebox}{一阶逻辑的组成}
一阶逻辑比命题逻辑多了：
\begin{itemize}
  \item \kw{个体域} $D$：讨论对象的集合。
  \item \kw{个体词}：表示对象，如 $a,b,x,y$。
  \item \kw{谓词}：表示对象性质或关系，如 $P(x),R(x,y)$。
  \item \kw{量词}：全称量词 $\forall$，存在量词 $\exists$。
\end{itemize}
\end{corebox}

\begin{methodbox}{一阶逻辑符号化常见句式}
\begin{center}
\begin{tabularx}{\textwidth}{>{\bfseries}p{4.2cm}X}
\toprule
自然语言 & 一阶逻辑形式 \\
\midrule
所有 $A$ 都是 $B$ & $\forall x(A(x)\to B(x))$ \\
有些 $A$ 是 $B$ & $\exists x(A(x)\land B(x))$ \\
没有 $A$ 是 $B$ & $\neg\exists x(A(x)\land B(x))$，等价于 $\forall x(A(x)\to\neg B(x))$ \\
有些 $A$ 不是 $B$ & $\exists x(A(x)\land\neg B(x))$ \\
恰有一个 $x$ 满足 $P(x)$ & $\exists x(P(x)\land\forall y(P(y)\to y=x))$ \\
\bottomrule
\end{tabularx}
\end{center}
\end{methodbox}

\subsection{一阶逻辑等值式与前束范式（对应教材 2.7）}

\begin{corebox}{量词否定等值式}
\[
\neg\forall x A(x)\equiv \exists x\neg A(x),\qquad
\neg\exists x A(x)\equiv \forall x\neg A(x).
\]
记忆：\kw{否定穿过量词时，量词互换，内部取反}。
\end{corebox}

\begin{corebox}{量词辖域收缩与扩张等值式}
若 $B$ 中不含 $x$ 的自由出现，则：
\begin{align*}
\forall x(A(x)\lor B)&\equiv \forall xA(x)\lor B,\\
\forall x(A(x)\land B)&\equiv \forall xA(x)\land B,\\
\forall x(A(x)\to B)&\equiv \exists xA(x)\to B,\\
\forall x(B\to A(x))&\equiv B\to\forall xA(x),\\[0.4em]
\exists x(A(x)\lor B)&\equiv \exists xA(x)\lor B,\\
\exists x(A(x)\land B)&\equiv \exists xA(x)\land B,\\
\exists x(A(x)\to B)&\equiv \forall xA(x)\to B,\\
\exists x(B\to A(x))&\equiv B\to\exists xA(x).
\end{align*}
\end{corebox}

\begin{corebox}{量词分配与量词交换}
\begin{align*}
\forall x(A(x)\land B(x))&\equiv \forall xA(x)\land\forall xB(x),\\
\exists x(A(x)\lor B(x))&\equiv \exists xA(x)\lor\exists xB(x),\\
\forall x\forall y A(x,y)&\equiv\forall y\forall x A(x,y),\\
\exists x\exists y A(x,y)&\equiv\exists y\exists x A(x,y).
\end{align*}
注意：$\forall$ 对 $\lor$、$\exists$ 对 $\land$ 一般不能分配；$\forall x\exists y$ 与 $\exists y\forall x$ 一般不能交换。
\end{corebox}

\begin{methodbox}{前束范式计算步骤}
前束范式形如
\[
Q_1x_1Q_2x_2\cdots Q_kx_k B,
\]
其中 $Q_i\in\{\forall,\exists\}$，$B$ 不含量词。
\begin{enumerate}
  \item 消去 $\to,\leftrightarrow$。
  \item 否定号内移，只作用到原子公式。
  \item 必要时改名，避免变量冲突。
  \item 用辖域扩张/收缩把所有量词移到公式最前面。
\end{enumerate}
\end{methodbox}

\section{集合、关系与函数部分（34 分）}

\subsection{集合运算与基本公式}

\begin{corebox}{集合基本运算}
设全集为 $E$，集合为 $A,B,C$。
\[
\begin{aligned}
A\cup B&=\set{x\mid x\in A\lor x\in B},\\
A\cap B&=\set{x\mid x\in A\land x\in B},\\
A-B&=\set{x\mid x\in A\land x\notin B}=A\cap\overline B,\\
\overline A&=E-A,\\
A\oplus B&=(A-B)\cup(B-A).
\end{aligned}
\]
\end{corebox}

\begin{corebox}{集合运算律速查}
\begin{longtable}{p{2.6cm}p{11cm}}
\toprule
律 & 公式 \\
\midrule
幂等律 & $A\cup A=A,\quad A\cap A=A$ \\
结合律 & $(A\cup B)\cup C=A\cup(B\cup C)$；$(A\cap B)\cap C=A\cap(B\cap C)$ \\
交换律 & $A\cup B=B\cup A$；$A\cap B=B\cap A$ \\
分配律 & $A\cup(B\cap C)=(A\cup B)\cap(A\cup C)$；$A\cap(B\cup C)=(A\cap B)\cup(A\cap C)$ \\
同一律 & $A\cup\varnothing=A$；$A\cap E=A$ \\
零律 & $A\cup E=E$；$A\cap\varnothing=\varnothing$ \\
排中律 & $A\cup\overline A=E$ \\
矛盾律 & $A\cap\overline A=\varnothing$ \\
吸收律 & $A\cup(A\cap B)=A$；$A\cap(A\cup B)=A$ \\
德摩根律 & $\overline{A\cup B}=\overline A\cap\overline B$；$\overline{A\cap B}=\overline A\cup\overline B$ \\
双重否定 & $\overline{\overline A}=A$ \\
差集公式 & $A-B=A\cap\overline B$；$A-(B\cup C)=(A-B)\cap(A-C)$；$A-(B\cap C)=(A-B)\cup(A-C)$ \\
对称差 & $A\oplus B=B\oplus A$；$(A\oplus B)\oplus C=A\oplus(B\oplus C)$；$A\oplus\varnothing=A$；$A\oplus A=\varnothing$ \\
\bottomrule
\end{longtable}
\end{corebox}

\subsection{证明集合包含或相等}

\begin{methodbox}{集合证明模板}
证明 $A\subseteq B$：取任意 $x$，假设 $x\in A$，经过等价变形推出 $x\in B$。\\
证明 $A=B$：证明 $A\subseteq B$ 且 $B\subseteq A$。\\
证明集合恒等式：可用元素法、运算律化简法、文氏图辅助法。
\end{methodbox}

\begin{exambox}{例：证明 $A-(B\cup C)=(A-B)\cap(A-C)$}
对任意 $x$，
\[
\begin{aligned}
x\in A-(B\cup C)
&\Leftrightarrow x\in A\land x\notin B\cup C\\
&\Leftrightarrow x\in A\land\neg(x\in B\lor x\in C)\\
&\Leftrightarrow x\in A\land x\notin B\land x\notin C\\
&\Leftrightarrow x\in A-B\land x\in A-C\\
&\Leftrightarrow x\in (A-B)\cap(A-C).
\end{aligned}
\]
\end{exambox}

\subsection{关系的定义与运算（对应教材 3.3）}

\begin{corebox}{关系的基本概念}
从 $A$ 到 $B$ 的二元关系是 $A\times B$ 的子集。若 $R\subseteq A\times B$，则
\[
\dom R=\set{x\mid \exists y(\langle x,y\rangle\in R)},\quad
\ran R=\set{y\mid \exists x(\langle x,y\rangle\in R)}.
\]
当 $R\subseteq A\times A$ 时，称 $R$ 为 $A$ 上的关系。
\end{corebox}

\begin{corebox}{关系运算}
设 $F,G$ 为关系，教材常用约定：$F\circ G$ 表示\kw{先经 $G$，再经 $F$}。
\begin{align*}
F^{-1}&=\set{\langle y,x\rangle\mid \langle x,y\rangle\in F},\\
F\circ G&=\set{\langle x,y\rangle\mid \exists z(\langle x,z\rangle\in G\land\langle z,y\rangle\in F)},\\
F\upharpoonright A&=\set{\langle x,y\rangle\mid \langle x,y\rangle\in F\land x\in A},\\
F[A]&=\ran(F\upharpoonright A).
\end{align*}
\end{corebox}

\begin{corebox}{关系运算性质}
\begin{align*}
(F^{-1})^{-1}&=F,\\
\dom F^{-1}&=\ran F,\\ \ran F^{-1}=\dom F,\\
(F\circ G)\circ H&=F\circ(G\circ H),\\
(F\circ G)^{-1}&=G^{-1}\circ F^{-1},\\
F\circ(G\cup H)&=(F\circ G)\cup(F\circ H),\\
(G\cup H)\circ F&=(G\circ F)\cup(H\circ F).
\end{align*}
注意：复合对交集一般只有包含关系，不可随意写成等式。
\end{corebox}

\begin{corebox}{关系的幂}
设 $R$ 为 $A$ 上的关系：
\[
R^0=I_A=\set{\langle x,x\rangle\mid x\in A},\qquad R^n=R^{n-1}\circ R\quad(n\ge1).
\]
常用性质：
\[
R^m\circ R^n=R^{m+n},\qquad (R^m)^n=R^{mn}.
\]
图上理解：$\langle x,y\rangle\in R^k$ 表示从 $x$ 到 $y$ 存在长度为 $k$ 的通路。
\end{corebox}

\subsection{关系的性质与闭包（对应教材 3.4）}

\begin{corebox}{五种性质}
设 $R$ 是 $A$ 上的关系。
\begin{center}
\begin{tabularx}{\textwidth}{p{2.5cm}X}
\toprule
性质 & 判定 \\
\midrule
自反性 & $\forall x\in A,\langle x,x\rangle\in R$ \\
反自反性 & $\forall x\in A,\langle x,x\rangle\notin R$ \\
对称性 & $\langle x,y\rangle\in R\Rightarrow\langle y,x\rangle\in R$ \\
反对称性 & $\langle x,y\rangle,\langle y,x\rangle\in R\Rightarrow x=y$ \\
传递性 & $\langle x,y\rangle,\langle y,z\rangle\in R\Rightarrow\langle x,z\rangle\in R$ \\
\bottomrule
\end{tabularx}
\end{center}
\end{corebox}

\begin{corebox}{关系运算保持性质表}
表中 $\checkmark$ 表示：若原关系具有该性质，则运算结果必然保持该性质。
\begin{center}
\begin{tabular}{c|ccccc}
\toprule
运算 & 自反 & 反自反 & 对称 & 反对称 & 传递 \\
\midrule
$R^{-1}$ & $\checkmark$ & $\checkmark$ & $\checkmark$ & $\checkmark$ & $\checkmark$ \\
$R_1\cap R_2$ & $\checkmark$ & $\checkmark$ & $\checkmark$ & $\checkmark$ & $\checkmark$ \\
$R_1\cup R_2$ & $\checkmark$ & $\checkmark$ & $\checkmark$ & $\times$ & $\times$ \\
$R_1-R_2$ & $\times$ & $\checkmark$ & $\checkmark$ & $\checkmark$ & $\times$ \\
$R_1\circ R_2$ & $\checkmark$ & $\times$ & $\times$ & $\times$ & $\times$ \\
\bottomrule
\end{tabular}
\end{center}
\end{corebox}

\begin{corebox}{闭包}
闭包是在原关系基础上添加尽量少的有序对，使其满足指定性质。
\[
r(R)=R\cup I_A,
\]
\[
s(R)=R\cup R^{-1},
\]
若 $A$ 有 $n$ 个元素，则传递闭包可写为
\[
t(R)=R\cup R^2\cup\cdots\cup R^n.
\]
\end{corebox}

\subsection{等价关系、等价类与划分（对应教材 3.6）}

\begin{corebox}{等价关系}
设 $R$ 是非空集合 $A$ 上的关系。若 $R$ 同时满足\kw{自反性、对称性、传递性}，则称 $R$ 为 $A$ 上的等价关系。
\end{corebox}

\begin{corebox}{等价类与商集}
若 $R$ 是 $A$ 上的等价关系，对 $x\in A$，
\[
[x]_R=\set{y\mid y\in A\land xRy}
\]
称为 $x$ 关于 $R$ 的等价类。所有等价类构成的集合称为商集：
\[
A/R=\set{[x]_R\mid x\in A}.
\]
\end{corebox}

\begin{corebox}{等价类重要性质}
\[
[x]_R\neq\varnothing,\\quad x\in[x]_R,
\]
\[
xRy\Rightarrow[x]_R=[y]_R,
\]
\[
\neg(xRy)\Rightarrow[x]_R\cap[y]_R=\varnothing,
\]
\[
\bigcup_{x\in A}[x]_R=A.
\]
\end{corebox}

\begin{corebox}{划分与等价关系一一对应}
集合 $A$ 的一个划分 $\pi$ 是 $A$ 的非空子集族，满足：
\begin{enumerate}
  \item 每个块非空；
  \item 任意两个不同块不相交；
  \item 所有块的并集为 $A$。
\end{enumerate}
一个等价关系的所有等价类构成 $A$ 的划分；反过来，一个划分也能确定一个等价关系。
\end{corebox}

\subsection{偏序关系与哈斯图（对应教材 3.7）}

\begin{corebox}{偏序关系}
设 $R$ 是非空集合 $A$ 上的关系。若 $R$ 同时满足\kw{自反性、反对称性、传递性}，则称 $R$ 为 $A$ 上的偏序关系，记为
\[
\langle A,R\rangle\quad\text{或}\quad \langle A,\preceq\rangle.
\]
若任意 $x,y\in A$ 都可比，则称为全序关系。
\end{corebox}

\begin{corebox}{偏序中的常考概念}
设 $\langle A,\preceq\rangle$ 为偏序集，$B\subseteq A$。
\begin{itemize}
  \item $x,y$ 可比：$x\preceq y$ 或 $y\preceq x$。
  \item $y$ 盖住 $x$：$x\prec y$，且不存在 $z$ 使 $x\prec z\prec y$。
  \item 极小元：$B$ 中没有比它更小的元素；最小元：它小于等于 $B$ 中所有元素。
  \item 极大元：$B$ 中没有比它更大的元素；最大元：它大于等于 $B$ 中所有元素。
  \item 下界：小于等于 $B$ 中所有元素的元素；上界：大于等于 $B$ 中所有元素的元素。
  \item 下确界：最大下界；上确界：最小上界。
\end{itemize}
\end{corebox}

\begin{methodbox}{哈斯图画法}
\begin{enumerate}
  \item 先确认关系是偏序关系。
  \item 删除所有自环。
  \item 删除由传递性推出的边，只保留覆盖关系。
  \item 较小元素画在下方，较大元素画在上方，箭头通常省略。
\end{enumerate}
\end{methodbox}

\subsection{函数的定义与性质（对应教材 3.8）}

\begin{corebox}{函数作为特殊关系}
函数是特殊的二元关系。关系 $f$ 为函数，当且仅当对任意 $x\in\dom f$，存在唯一的 $y\in\ran f$ 使 $\langle x,y\rangle\in f$。
\end{corebox}

\begin{corebox}{从 $A$ 到 $B$ 的函数写法}
设 $A,B$ 为集合，若函数 $f$ 满足
\[
\dom f=A,
\qquad
\ran f\subseteq B,
\]
则称 $f$ 是从 $A$ 到 $B$ 的函数，记作
\[
f:A\to B.
\]
从 $A$ 到 $B$ 的所有函数构成集合
\[
B^A=\set{f\mid f:A\to B}.
\]
若 $A,B$ 为有限集，则 $|B^A|=|B|^{|A|}$。
\end{corebox}

\begin{corebox}{单射、满射、双射}
\begin{itemize}
  \item 单射：$x_1\neq x_2\Rightarrow f(x_1)\neq f(x_2)$。
  \item 满射：$\ran f=B$。
  \item 双射：既单射又满射。只有双射才有真正意义上的反函数 $f^{-1}:B\to A$。
\end{itemize}
\end{corebox}

\begin{corebox}{特殊函数}
\begin{itemize}
  \item 特征函数：对 $A\subseteq E$，
  \[
  \chi_A(x)=\begin{cases}1,&x\in A,\\0,&x\notin A.
  \end{cases}
  \]
  \item 自然映射：若 $R$ 是 $A$ 上的等价关系，定义
  \[
  \eta:A\to A/R,\qquad \eta(x)=[x]_R.
  \]
\end{itemize}
\end{corebox}

\section{图的基本概念与矩阵表示（图论基础）}

\subsection{图、度数与握手定理}

\begin{corebox}{图的基本定义}
无向图或有向图可记为
\[
G=\langle V,E\rangle.
\]
有 $n$ 个顶点的图称为 $n$ 阶图。没有边的图称为零图；一阶零图只有一个顶点，也称平凡图。
\end{corebox}

\begin{corebox}{端点、关联、相邻、环、孤立点}
无向边 $e=(u,v)$ 中，$u,v$ 是 $e$ 的端点，边 $e$ 与 $u,v$ 关联。若两个顶点由同一条边连接，则称相邻。端点重合的边称为环。不与任何边关联的顶点称为孤立点。
\end{corebox}

\begin{corebox}{度数}
无向图中，顶点 $v$ 的度数记为 $d(v)$。环对度数贡献 $2$。有向图中：
\[
d^+(v)=\text{出度},\qquad d^-(v)=\text{入度},\qquad d(v)=d^+(v)+d^-(v).
\]
最大度、最小度：
\[
\Delta(G)=\max\set{d(v)\mid v\in V},\qquad
\delta(G)=\min\set{d(v)\mid v\in V}.
\]
有向图还可定义最大出度、最小出度、最大入度、最小入度：
\[
\Delta^+(D),\delta^+(D),\Delta^-(D),\delta^-(D).
\]
\end{corebox}

\begin{corebox}{握手定理}
设图 $G=\langle V,E\rangle$，$V=\set{v_1,\ldots,v_n}$，$|E|=m$，则
\[
\sum_{i=1}^n d(v_i)=2m.
\]
推论：任意图中，\kw{奇度顶点的个数为偶数}。
\end{corebox}

\begin{corebox}{有向图度数定理}
设有向图 $D=\langle V,E\rangle$，$V=\set{v_1,\ldots,v_n}$，$|E|=m$，则
\[
\sum_{i=1}^n d^+(v_i)=\sum_{i=1}^n d^-(v_i)=m.
\]
\end{corebox}

\subsection{简单图、完全图、子图、补图与同构}

\begin{corebox}{平行边、多重图、简单图}
无向图中，关联同一对顶点的无向边若多于一条，称为平行边。含平行边的图称为多重图。不含环也不含平行边的图称为简单图。
\end{corebox}

\begin{corebox}{完全图与子图}
$n$ 阶无向完全图记为 $K_n$，任意两个不同顶点都相邻。若 $G'=(V',E')$ 满足 $V'\subseteq V,E'\subseteq E$，则 $G'$ 是 $G$ 的子图。若 $V'=V$，则为生成子图。由顶点子集诱导的子图记为 $G[V']$。
\end{corebox}

\begin{corebox}{补图与同构}
对 $n$ 阶无向简单图 $G$，其补图 $\overline G$ 与 $G$ 顶点集相同，且两顶点在 $\overline G$ 中相邻当且仅当它们在 $G$ 中不相邻。\\
图同构：若存在双射 $f:V_1\to V_2$，使得任意 $u,v\in V_1$，$uv\in E_1\Leftrightarrow f(u)f(v)\in E_2$，则 $G_1\cong G_2$。
\end{corebox}

\subsection{图的矩阵表示}

\begin{corebox}{无向图关联矩阵}
设无向图 $G=\langle V,E\rangle$，$V=\set{v_1,\ldots,v_n}$，$E=\set{e_1,\ldots,e_m}$。关联矩阵 $M(G)=(m_{ij})_{n\times m}$，其中 $m_{ij}$ 表示顶点 $v_i$ 与边 $e_j$ 的关联次数。
\begin{itemize}
  \item 每列元素和为 $2$：普通边对应两个 $1$，环可记为一个 $2$。
  \item 第 $i$ 行元素和为 $d(v_i)$。
  \item 所有元素之和为 $2m$。
\end{itemize}
\end{corebox}

\begin{corebox}{有向图关联矩阵}
设有向图无环，若 $v_i$ 是边 $e_j$ 的始点，记 $1$；若 $v_i$ 是边 $e_j$ 的终点，记 $-1$；不关联记 $0$。
\begin{itemize}
  \item 每列恰有一个 $1$ 和一个 $-1$。
  \item 第 $i$ 行中 $1$ 的个数等于 $d^+(v_i)$，$-1$ 的个数等于 $d^-(v_i)$。
\end{itemize}
\end{corebox}

\begin{corebox}{有向图邻接矩阵与矩阵幂}
邻接矩阵 $A(D)=(a_{ij})$ 中，$a_{ij}$ 表示从 $v_i$ 到 $v_j$ 的边数。
\begin{itemize}
  \item 第 $i$ 行元素和为 $d^+(v_i)$。
  \item 第 $j$ 列元素和为 $d^-(v_j)$。
  \item 所有元素之和为 $m$。
  \item $A^l$ 中第 $(i,j)$ 个元素表示从 $v_i$ 到 $v_j$ 的长度为 $l$ 的通路条数；对角线元素表示回路条数。
\end{itemize}
若 $B_r=A+A^2+\cdots+A^r$，则 $B_r$ 的元素可反映长度不超过 $r$ 的通路总数。
\end{corebox}

\begin{corebox}{可达矩阵}
可达矩阵 $P(D)=(p_{ij})$：
\[
p_{ij}=\begin{cases}
1,&v_i\text{ 可达 }v_j,\\
0,&\text{否则}.
\end{cases}
\]
由于每个顶点到自身可达，所以可达矩阵主对角线元素为 $1$。
\end{corebox}

\subsection{连通性、割点与割边}

\begin{corebox}{无向图连通性}
无向图中，若任意两个顶点之间都有路径，则称图连通。删除某个顶点及其关联边后使连通分支数增加，则该顶点为割点。删除某条边后使连通分支数增加，则该边为桥或割边。
\end{corebox}

\begin{corebox}{有向图连通性}
\begin{itemize}
  \item 强连通：任意两个顶点互相可达。
  \item 单向连通：任意两个顶点中，至少一个能到达另一个。
  \item 弱连通：忽略方向后得到的无向图连通。
\end{itemize}
判定顺序：强连通 $\Rightarrow$ 单向连通 $\Rightarrow$ 弱连通。
\end{corebox}

\begin{methodbox}{点割集与边割集}
点割集：删除某个顶点子集后，图不连通，且该子集通常要求极小。边割集类似，删除边集后图不连通。考试中按题目要求列出所有最小点割集或边割集时，要检查“删前连通、删后不连通、任何真子集删后仍连通”。
\end{methodbox}

\subsection{最短路径与图着色}

\begin{methodbox}{Dijkstra 最短路径算法}
适用于边权非负的图。
\begin{enumerate}
  \item 初始：源点距离为 $0$，其余为 $\infty$。
  \item 每次从未确定顶点中选距离最小者，标记为确定。
  \item 用该顶点松弛它的邻边。
  \item 重复直到目标顶点确定或所有顶点确定。
\end{enumerate}
\end{methodbox}

\begin{corebox}{图着色}
对无向图顶点着色，使相邻顶点颜色不同。最少颜色数称为色数，记为 $\chi(G)$。常用结论：
\begin{itemize}
  \item 空图色数为 $1$。
  \item 完全图 $K_n$ 的色数为 $n$。
  \item 非空二部图的色数为 $2$。
  \item 奇圈 $C_{2k+1}$ 色数为 $3$，偶圈 $C_{2k}$ 色数为 $2$。
\end{itemize}
\end{corebox}

\section{特殊的图}

\subsection{二部图}

\begin{corebox}{二部图与完全二部图}
无向图 $G=\langle V,E\rangle$ 是二部图，若 $V$ 可划分为两个互不相交的非空子集 $V_1,V_2$，使得每条边的两个端点一个属于 $V_1$，另一个属于 $V_2$。记为
\[
G=\langle V_1,V_2,E\rangle.
\]
若 $V_1$ 中任意顶点都与 $V_2$ 中任意顶点相邻，则为完全二部图，记为 $K_{m,n}$。
\end{corebox}

\begin{corebox}{二部图判定定理}
无向图 $G$ 是二部图，当且仅当 $G$ 中没有奇长度回路。
\end{corebox}

\begin{corebox}{匹配、最大匹配、完美匹配}
设 $M\subseteq E$。若 $M$ 中任意两条边均不相邻，则称 $M$ 为匹配。边数最多的匹配称为最大匹配。若图中每个顶点都是 $M$ 中某条边的端点，则称 $M$ 为完美匹配。
\end{corebox}

\begin{corebox}{二部图中的完全匹配与 Hall 定理}
设二部图 $G=\langle V_1,V_2,E\rangle$，若匹配 $M$ 中每条边均连接 $V_1$ 与 $V_2$，且 $V_1$ 中每个顶点都被匹配，则称 $M$ 为 $G$ 中从 $V_1$ 到 $V_2$ 的完全匹配。少的一方匹配完即可。\\
Hall 定理：若 $|V_1|\le |V_2|$，则 $G$ 存在从 $V_1$ 到 $V_2$ 的完全匹配，当且仅当 $V_1$ 中任意 $k$ 个顶点至少邻接 $V_2$ 中的 $k$ 个顶点。
\end{corebox}

\subsection{欧拉图}

\begin{corebox}{欧拉通路、欧拉回路与欧拉图}
经过图中每条边一次且仅一次的通路称为欧拉通路；若它是回路，则称为欧拉回路。存在欧拉回路的图称为欧拉图。存在欧拉通路但不存在欧拉回路的图称为半欧拉图。
\end{corebox}

\begin{corebox}{无向图欧拉判定}
设无向图 $G$ 连通（孤立点通常不影响非零边所在部分的判断）：
\begin{itemize}
  \item $G$ 有欧拉回路 $\Leftrightarrow$ 所有顶点度数均为偶数。
  \item $G$ 有欧拉通路但无欧拉回路 $\Leftrightarrow$ 恰有两个奇度顶点。
\end{itemize}
\end{corebox}

\begin{corebox}{有向图欧拉判定}
设有向图 $D$ 连通：
\begin{itemize}
  \item $D$ 有欧拉回路 $\Leftrightarrow$ 每个顶点入度等于出度。
  \item $D$ 有欧拉通路 $\Leftrightarrow$ 除两个特殊顶点外，其余顶点入度等于出度；起点满足 $d^+(v)=d^-(v)+1$，终点满足 $d^-(v)=d^+(v)+1$。
\end{itemize}
\end{corebox}

\subsection{哈密顿图}

\begin{corebox}{哈密顿通路、哈密顿回路与哈密顿图}
经过图中每个顶点一次且仅一次的通路称为哈密顿通路；若它是回路，则称为哈密顿回路。存在哈密顿回路的图称为哈密顿图；存在哈密顿通路的图常称为半哈密顿图。
\end{corebox}

\begin{corebox}{哈密顿图必要条件}
若无向图 $G=\langle V,E\rangle$ 是哈密顿图，则对任意非空真子集 $S\subset V$，有
\[
p(G-S)\le |S|,
\]
其中 $p(G-S)$ 表示删除 $S$ 后所得图的连通分支数。若找到某个 $S$ 使 $p(G-S)>|S|$，即可判定 $G$ 不是哈密顿图。
\end{corebox}

\begin{corebox}{哈密顿通路/回路充分条件}
设 $G$ 是 $n(n\ge3)$ 阶无向简单图。
\begin{itemize}
  \item 若任意两个不相邻顶点 $u,v$ 都满足 $d(u)+d(v)\ge n-1$，则 $G$ 存在哈密顿通路。
  \item 若任意两个不相邻顶点 $u,v$ 都满足 $d(u)+d(v)\ge n$，则 $G$ 是哈密顿图。
\end{itemize}
这些是充分条件，不满足时不能直接推出“不是”。
\end{corebox}

\section{树}

\subsection{无向树及生成树}

\begin{corebox}{树、森林、树叶、分支点}
不含回路的连通无向图称为无向树，简称树。每个连通分支都是树的非连通无向图称为森林。平凡图称为平凡树。树中度数为 $1$ 的顶点称为树叶，度数大于等于 $2$ 的顶点称为分支点。
\end{corebox}

\begin{corebox}{树的等价命题}
设 $G=\langle V,E\rangle$，$|V|=n$，$|E|=m$。下列命题等价：
\begin{enumerate}
  \item $G$ 是无向树；
  \item $G$ 的任意两个顶点之间有唯一一条路径；
  \item $G$ 连通且 $m=n-1$；
  \item $G$ 无回路且 $m=n-1$；
  \item $G$ 无回路，但在任意两个不相邻顶点之间加一条边，就形成唯一一条初级回路；
  \item $G$ 连通，且每条边都是桥。
\end{enumerate}
核心公式：
\[
\boxed{\text{树的边数}=\text{顶点数}-1}\qquad m=n-1.
\]
\end{corebox}

\begin{corebox}{树叶结论与生成树}
$n$ 阶非平凡树中至少有两片树叶。\\
设 $G$ 是无向连通图。若 $T$ 是 $G$ 的生成子图且 $T$ 是树，则称 $T$ 是 $G$ 的生成树。生成树中的边称为树枝，不在生成树中的边称为弦，所有弦构成余树。
\end{corebox}

\begin{corebox}{生成树与基本回路}
任意无向连通图都有生成树。若无向连通图有 $n$ 个顶点、$m$ 条边，则 $m\ge n-1$。\\
设 $T$ 是连通图 $G$ 的一棵生成树。对每条弦 $e$，把 $e$ 加入 $T$ 会形成唯一一个初级回路，称为对应于弦 $e$ 的基本回路。
\end{corebox}

\subsection{根树及其应用}

\begin{corebox}{有向树与根树}
一个有向图 $D$，若略去有向边方向后得到的无向图是一棵无向树，则称 $D$ 为有向树。\\
非平凡有向树中，入度为 $0$ 的顶点称为根；入度为 $1$、出度为 $0$ 的顶点称为树叶；入度为 $1$、出度大于 $0$ 的顶点称为内点或分支点。若有且仅有一个顶点入度为 $0$，其余顶点入度均为 $1$，则称为根树。
\end{corebox}

\begin{corebox}{根树中的亲属关系、层数与高度}
若有边 $a\to b$，则 $a$ 是 $b$ 的父亲，$b$ 是 $a$ 的儿子。有同一个父亲的顶点互为兄弟。若从 $a$ 可达 $b$，则 $a$ 是 $b$ 的祖先，$b$ 是 $a$ 的后代。\\
从根到顶点 $v$ 的通路长度称为 $v$ 的层数；树中顶点最大层数称为树的高度，记为 $h(T)$。
\end{corebox}

\begin{corebox}{有序树、$r$ 叉树与二叉树}
若根树每层顶点规定次序，则称为有序树。若每个分支点至多有 $r$ 个儿子，则称为 $r$ 叉树。若每个分支点恰有 $r$ 个儿子，则称为正则 $r$ 叉树。若所有树叶层数相同，则称为完全正则 $r$ 叉树。$r=2$ 时为二叉树。
\end{corebox}

\subsection{最小生成树：Kruskal 算法}

\begin{corebox}{最小生成树}
带权连通无向图的生成树中，边权总和最小者称为最小生成树。
\end{corebox}

\begin{methodbox}{Kruskal 算法}
\begin{enumerate}
  \item 按边权从小到大排序。
  \item 依次考察每条边：若加入后不形成回路，则加入；否则跳过。
  \item 直到选出 $n-1$ 条边为止。
\end{enumerate}
关键检查：\kw{不能形成回路}。
\end{methodbox}

\subsection{哈夫曼树与哈夫曼编码}

\begin{corebox}{带权路径长度 WPL}
若树叶 $v_i$ 的权值为 $w_i$，从根到 $v_i$ 的路径长度为 $l_i$，则
\[
WPL=\sum_i w_i l_i.
\]
哈夫曼树是 WPL 最小的二叉树，用于构造最优前缀编码。
\end{corebox}

\begin{methodbox}{哈夫曼算法}
\begin{enumerate}
  \item 将每个权值看成一棵单节点树。
  \item 每次选权值最小的两棵树合并，新树权值为二者之和。
  \item 重复直到只剩一棵树。
  \item 从根开始，左分支可记 $0$，右分支可记 $1$，读出各叶子的编码。
\end{enumerate}
\end{methodbox}

\subsection{二叉树遍历与波兰/逆波兰表达式}

\begin{corebox}{二叉树遍历}
\begin{itemize}
  \item 前序遍历：根 - 左 - 右。
  \item 中序遍历：左 - 根 - 右。
  \item 后序遍历：左 - 右 - 根。
  \item 层序遍历：从上到下、从左到右。
\end{itemize}
已知前序 + 中序，或后序 + 中序，可以唯一确定二叉树。只有前序 + 后序一般不能唯一确定，除非题目另有限制。
\end{corebox}

\begin{corebox}{表达式树与波兰符号}
表达式树中，运算符为内部结点，操作数为叶子。
\begin{itemize}
  \item 前序遍历得到前缀表达式，即波兰符号。
  \item 中序遍历得到中缀表达式，需要括号辅助。
  \item 后序遍历得到后缀表达式，即逆波兰符号。
\end{itemize}
\end{corebox}

\section{考前速查模板}

\begin{exambox}{逻辑题模板}
\begin{itemize}
  \item 符号化：先定义命题/谓词，再写公式。
  \item 判断公式类型：真值表或等值变形。
  \item 主范式：按变元顺序列真值表，真行写小项，假行写大项。
  \item 推理证明：每行公式 + 每行依据，最后推出目标结论。
  \item 前束范式：消蕴含、否定内移、变量改名、量词前移。
\end{itemize}
\end{exambox}

\begin{exambox}{关系题模板}
\begin{itemize}
  \item 判断性质：先查自环，再查反向边，再查传递三元组。
  \item 等价关系：自反 + 对称 + 传递；再写等价类 $[x]_R$ 和商集 $A/R$。
  \item 偏序关系：自反 + 反对称 + 传递；画 Hasse 图时删自环、删传递边。
  \item 闭包：自反闭包加 $I_A$；对称闭包加逆关系；传递闭包加 $R^2,R^3,\ldots$。
\end{itemize}
\end{exambox}

\begin{exambox}{图论题模板}
\begin{itemize}
  \item 度数/边数：优先用握手定理 $\sum d(v)=2m$。
  \item 有向图：出度和 = 入度和 = 边数。
  \item 矩阵幂：$A^l$ 的 $(i,j)$ 元素表示长度为 $l$ 的通路数。
  \item 二部图：看是否有奇回路；匹配题看 Hall 条件。
  \item 欧拉图：看边是否能一次走完，核心查度数。
  \item 哈密顿图：看顶点是否能一次经过，不能用欧拉条件替代。
  \item 树：看到树先写 $m=n-1$；生成树有 $n-1$ 条边。
  \item Kruskal：按权从小到大选，不成环才加入。
\end{itemize}
\end{exambox}


\section{模拟卷查缺补漏与冲刺计划}

\begin{corebox}{按两套模拟卷补强的题型结构}
\begin{center}
\begin{tabular}{p{3.4cm}p{2.2cm}p{2.4cm}p{5.5cm}}
\toprule
题型 & 题量 & 总分 & 主要考法 \\
\midrule
单选题 & 15 道 & 30 分 & 逻辑、集合、关系、函数、图论基础概念 \\
判断题 & 10 道 & 10 分 & 定义正误、性质判定、图论结论 \\
计算证明题 & 5 道 & 30 分 & 等值演算、主范式、关系、函数、图判定 \\
综合应用题 & 3 道 & 30 分 & 最短路、Kruskal、Huffman、树与矩阵 \\
\bottomrule
\end{tabular}
\end{center}
\end{corebox}

\begin{pitfallbox}{查缺补漏重点}
\begin{enumerate}
  \item 等值演算证明必须逐步写依据，不能只写首尾等价。
  \item 主析取范式与主合取范式要按变元固定顺序写小项/大项编号。
  \item 关系性质题不能只凭直觉，应逐项检查自反、对称、反对称、传递。
  \item 函数题要分清单射、满射、双射，尤其注意陪域是否完全覆盖。
  \item 图论综合题要写过程：Dijkstra/Kruskal/Huffman 不能只给最终结果。
\end{enumerate}
\end{pitfallbox}

\begin{methodbox}{7 天复习安排}
\begin{enumerate}
  \item Day 1：命题逻辑符号化、真值表、等值式。
  \item Day 2：主范式、构造证明、一阶逻辑符号化与前束范式。
  \item Day 3：集合运算证明、关系运算、关系性质与闭包。
  \item Day 4：等价关系、偏序关系、函数性质与复合反函数。
  \item Day 5：图的度数、连通性、矩阵表示、可达矩阵和最短路。
  \item Day 6：二部图、欧拉图、哈密顿图、树、Kruskal、Huffman。
  \item Day 7：做一套模拟卷，按错题回到对应模板复盘。
\end{enumerate}
\end{methodbox}

\begin{methodbox}{3 天冲刺安排}
\begin{enumerate}
  \item Day 1：背逻辑等值式、范式、关系性质和函数判定。
  \item Day 2：刷计算证明题，重点练主范式、闭包、等价类、图矩阵。
  \item Day 3：练综合应用题，背最后 30 分钟速记清单。
\end{enumerate}
\end{methodbox}

\begin{exambox}{最后 30 分钟速记清单}
\begin{itemize}
  \item $p\to q\equiv \neg p\lor q$；$\neg(p\land q)\equiv\neg p\lor\neg q$；$\neg(p\lor q)\equiv\neg p\land\neg q$。
  \item 主析取范式看真值为真的行；主合取范式看真值为假的行。
  \item 等价关系：自反、对称、传递；偏序关系：自反、反对称、传递。
  \item 自反闭包加 $I_A$；对称闭包加 $R^{-1}$；传递闭包加 $R^2,R^3,\ldots$ 或 Warshall。
  \item 单射看不同原像是否有不同像；满射看陪域每个元素是否被取到。
  \item 无向图握手定理：$\sum d(v)=2m$；树：$m=n-1$ 且连通无回路。
  \item 欧拉回路：连通且所有顶点度为偶数；欧拉通路：恰有 0 或 2 个奇度顶点。
  \item 二部图无奇回路；Kruskal 按边权从小到大选且不成环；Huffman 每次合并两个最小权。
\end{itemize}
\end{exambox}

\end{document}
